\chapter{Ceruloplasmin}
\section*{What Is This Test?}
Ceruloplasmin is an alpha-2 glycoprotein synthesized primarily by the liver that serves as the major copper-carrying protein in the blood, binding 6--7 copper atoms per molecule (90--95\% of serum copper is bound to ceruloplasmin). Ceruloplasmin functions as a ferroxidase, catalyzing the oxidation of ferrous iron (Fe$^{2+}$) to ferric iron (Fe$^{3+}$), which is required for the incorporation of iron into transferrin and the mobilization of iron from storage sites. Ceruloplasmin also has antioxidant properties (preventing free radical-mediated lipid peroxidation) and participates in copper transport to peripheral tissues. Serum ceruloplasmin measurement is most clinically useful for the diagnosis of Wilson disease (hepatolenticular degeneration) \cite{roberts2008diagnosis}, an autosomal recessive disorder of copper metabolism caused by mutations in the ATP7B gene (chromosome 13q14.3) encoding a copper-transporting ATPase. In Wilson disease, defective ATP7B function impairs the incorporation of copper into ceruloplasmin in the liver (producing apo-ceruloplasmin that is rapidly degraded, resulting in low serum ceruloplasmin) and impairs the biliary excretion of copper, leading to copper accumulation in the liver (chronic hepatitis, cirrhosis, acute liver failure), brain (basal ganglia, causing movement disorders --- dystonia, tremor, parkinsonism, dysarthria), and cornea (Kayser-Fleischer rings from copper deposition in Descemet membrane). Low ceruloplasmin (<20 mg/dL) is a hallmark of Wilson disease \cite{brewer2001wilsons}. However, ceruloplasmin is also an acute phase reactant and may be elevated in inflammation, infection, pregnancy, and estrogen therapy. Ceruloplasmin may also be low in heterozygotes for ATP7B mutations (carriers, asymptomatic), aceruloplasminemia, Menkes disease, protein-losing enteropathy/nephropathy, and malnutrition.
\section*{Alternative Names}
Ceruloplasmin; Serum Ceruloplasmin; Copper Oxidase; Ferroxidase; Cp
\section*{Principle}
Ceruloplasmin is measured by immunoturbidimetry or immunonephelometry using specific anti-ceruloplasmin antibodies. The ferroxidase enzymatic activity of ceruloplasmin can also be measured spectrophotometrically using p-phenylenediamine as a substrate (ceruloplasmin oxidase activity). Specimen: serum separator tube (SST, gold-top).
\section*{History}
Ceruloplasmin was discovered in 1948 by Holmberg and Laurell in Sweden, who isolated a blue copper-containing alpha-2 globulin from plasma and named it for its sky-blue (cerulean) color. In 1952, Scheinberg and Gitlin demonstrated that serum ceruloplasmin is markedly decreased in most patients with Wilson disease, providing the first broadly useful diagnostic test for the disorder. Holmberg and Laurell had earlier observed an oxidase activity, and in 1966 Osaki, Johnson, and Frieden established that ceruloplasmin functions as a ferroxidase, oxidizing ferrous iron (Fe$^{2+}$) to ferric iron (Fe$^{3+}$). The defective gene in Wilson disease, ATP7B, was mapped to chromosome 13 and cloned in 1993 by three independent groups. Clinical laboratory measurement of ceruloplasmin evolved from enzymatic oxidase assays to the immunoturbidimetric and immunonephelometric methods used today.
\section*{Why This Test Is Ordered}
\begin{itemize}
  \item Evaluation of suspected Wilson disease in patients with unexplained liver disease, cirrhosis or acute liver failure, movement disorders, psychiatric symptoms, or a positive family history
  \item Screening of asymptomatic siblings and family members of patients with Wilson disease
  \item Evaluation of unexplained Coombs-negative hemolytic anemia, which can accompany acute hepatic copper release in Wilson disease
  \item Assessment of copper metabolism in suspected aceruloplasminemia (low ceruloplasmin with anemia, diabetes, and neurodegeneration)
  \item Evaluation of suspected Menkes disease (X-linked copper deficiency) in infants
  \item Part of the diagnostic workup of unexplained low serum copper or suspected copper deficiency states
\end{itemize}
\section*{Primary vs Secondary Test}
Ceruloplasmin is a primary screening test for Wilson disease and a required component of the diagnostic score \cite{easl2012easl}, but it is not diagnostic on its own. Definitive diagnosis relies on confirmatory secondary testing, including 24-hour urinary copper, hepatic copper concentration, slit-lamp examination for Kayser-Fleischer rings, and ATP7B gene analysis.
\section*{How This Test Is Performed}
Blood is collected by standard venipuncture into a serum separator tube (SST, gold-top). Ceruloplasmin is measured in serum by immunoturbidimetry or immunonephelometry using specific anti-ceruloplasmin antibodies, which quantify the protein mass directly; the ferroxidase enzymatic activity can also be measured spectrophotometrically with p-phenylenediamine as substrate. The result is reported in milligrams per deciliter (mg/dL). No fasting is required, and results are typically available within 1--2 days.
\section*{What Preparation Is Needed}
No fasting or special preparation is required, and the sample may be drawn at any time of day. Because estrogen raises ceruloplasmin, the clinician should note pregnancy, oral contraceptive use, or hormone replacement therapy when interpreting the result.
\section*{Contraindications and Precautions}
\begin{itemize}
  \item Ceruloplasmin is an acute phase reactant: infection, inflammation, malignancy, pregnancy, and estrogen therapy elevate levels and can mask a genuinely low value in Wilson disease
  \item Up to 5--15\% of patients with Wilson disease, particularly those with active hepatic inflammation, have a normal ceruloplasmin level
  \item Heterozygous carriers of ATP7B mutations may have mildly low ceruloplasmin without copper accumulation and must not be labeled as affected
  \item Levels are age-dependent: neonates normally have lower ceruloplasmin (10--25 mg/dL)
  \item Units differ among laboratories (mg/dL versus mg/L); always verify the reference interval
  \item A low ceruloplasmin alone is not diagnostic of Wilson disease
\end{itemize}
\section*{Risks}
Minimal. Routine venipuncture risks include transient pain, bruising, and, rarely, hematoma, infection, or vasovagal syncope.
\section*{What Happens After the Test}
A low ceruloplasmin level triggers confirmatory evaluation: 24-hour urinary copper, slit-lamp examination for Kayser-Fleischer rings, hepatic copper quantification on liver biopsy, and ATP7B gene sequencing. A confirmed diagnosis of Wilson disease is treated lifelong with copper-chelating agents or zinc, and first-degree relatives should be screened. A normal ceruloplasmin in the appropriate clinical context does not exclude Wilson disease and may require further testing.
\section*{Understanding Results}
\subsection*{Decreased Ceruloplasmin (<20 mg/dL)}
Wilson disease: Low ceruloplasmin (<20 mg/dL) is found in >95\% of patients with Wilson disease \cite{ferenci2003diagnosis}. Homozygous or compound heterozygous ATP7B mutations. Acute liver failure from Wilson disease: markedly low ceruloplasmin, negative or normal alkaline phosphatase, positive Coombs-negative hemolytic anemia (from massive copper release from necrotic hepatocytes), and high 24-hour urinary copper. Aceruloplasminemia: Genetic absence of ceruloplasmin. Heterozygous ATP7B carriers: Mildly low ceruloplasmin (15--20 mg/dL) without copper accumulation. Menkes disease: Defective copper absorption (ATP7A copper transporter mutation, X-linked recessive). Malnutrition, protein-losing enteropathy/nephropathy, or end-stage liver disease (impaired hepatic synthesis).
\subsection*{Normal or Elevated Ceruloplasmin}
Does not exclude Wilson disease: up to 5--15\% of patients with Wilson disease have a normal ceruloplasmin level (particularly those with acute hepatic inflammation, which elevates ceruloplasmin as an acute phase reactant). Elevated ceruloplasmin is an acute phase reactant, elevated in: infection, inflammation, malignancy, pregnancy (elevated due to estrogen), estrogen therapy, and oral contraceptives.
\section*{Normal Results}
Serum ceruloplasmin: 20--60 mg/dL (adults) \cite{rifai2018tietz}. Neonates: lower (10--25 mg/dL). Wilson disease: <20 mg/dL (most patients), <10 mg/dL (severe Wilson disease). A low ceruloplasmin alone is NOT diagnostic of Wilson disease; the diagnosis is confirmed by: elevated 24-hour urinary copper (>100 mcg/24 hours, >40 mcg/24 hours in symptomatic patients), increased hepatic copper concentration on liver biopsy (>250 mcg/g dry weight), Kayser-Fleischer rings on slit-lamp examination (present in >95\% of patients with neurologic/psychiatric Wilson disease), and ATP7B gene mutation analysis.
\section*{Follow-up Tests}
24-hour urinary copper, hepatic copper concentration (liver biopsy), serum free copper (non-ceruloplasmin-bound copper, calculated as: serum copper --- [ceruloplasmin $\times$ 3.15], normal <10--15 mcg/dL), ATP7B gene sequencing, slit-lamp examination for Kayser-Fleischer rings, liver function tests, CBC (Coombs-negative hemolytic anemia in acute Wilson disease).
\section*{References}
\bibliographystyle{unsrtnat}
\bibliography{references}

\clearpage
